\setcounter{ex}{0}
\setcounter{paragraph}{0}
\PhanI
%\loai{Dùng công thức độc lập x, v tìm x}
\begin{ex}%[3P1B2-2][Loai1]
Một vật dao động điều hòa với biên độ A, vận tốc cực đại $v_{\max}$. Vật có tốc độ 0,6$v_{\max}$ khi li độ của vật có độ lớn là
\choice
{\True 0,8A}
{0,6A}
{0,4A}
{0,5A}
\loigiai{
\begin{itemize}
\item Luôn có: $\left(\dfrac{x}{A}\right)^2+\left(\dfrac{v}{v_{\max}}\right)^2=1\quad  (*)$
\item Khi $\left| v \right|=0,6v_{\max}\quad (*) \Rightarrow \left(\dfrac{x}{A}\right)^2+\left(\dfrac{0,6v_{\max}}{v_{\max}}\right)^2=1\Rightarrow \left|x\right|=0,8A$.
\end{itemize}
}

%\haidong{5}
\end{ex} \haidong{20}
%\loai{Dùng công thức độc lập x, v tìm A}
\begin{ex}%[3P1K2-2][Loai1]
Một chất điểm dao động điều hòa trên trục Ox. Trong thời gian 31,4 s chất điểm thực hiện được 100 dao động toàn phần. Lúc chất điểm đi qua vị trí có li độ 2 cm theo chiều âm với tốc độ là 40$\sqrt{3}$ cm/s. Lấy $\pi = 3,14$. Biên độ dao động của chất điểm là
\choice 
{\True 4 cm}
{2$\sqrt{2}$ cm}
{3 cm}
{2$\sqrt{3}$ cm}
\loigiai{
\begin{itemize}
\item Chu kì: $T = \dfrac{31,4}{100} = \dfrac{\pi}{10}\ s \Rightarrow \omega = 20\ rad/s$.
\item Sử dụng phương trình độc lập thời gian, ta được phương trình:
\begin{align*}
\dfrac{2^2}{A^2} + \dfrac{(40\sqrt{3})^2}{20^2.A^2}= 1\Rightarrow A = 4\ cm.
\end{align*}
\end{itemize}
}
%\haidong{4}
\end{ex} \haidong{20}
%\loai{Cho li độ x tìm v}
\begin{ex} %[3P1B2-2][Loai1] 
Vật dao động điều hoà với biên độ $A = 5$ cm, tần số $f = 4$ Hz. Tốc độ vật khi có li độ $x = 3$ cm là
\choice 
{$\left|v\right| = 2\pi$ cm/s}
{$\left|v\right|= 16\pi$ cm/s}
{\True $\left|v\right|= 32\pi$ cm/s}
{$\left|v\right|= 64\pi$ cm/s}
\loigiai{
\begin{itemize}
\item Ta có: $\omega = 2\pi f = 8\pi\ rad/s.$
\item Biểu thức độc lập: $A^2 = x^2 + \dfrac{v^2}{\omega^2}\Rightarrow |v| = \omega\sqrt{A^2-x^2} = 8\pi\sqrt{0,05^2-0,03^2} = 0,32\pi\ m/s = 32\pi\ cm/s$.
\end{itemize}
}
%\haidong{4}
\end{ex} \haidong{20}
%\loai{Cho khoảng cách từ vận tới vị trí cân bằng, tìm v}
\begin{ex}%[3P1B2-2][Loai1]
Một vật dao động điều hòa trên trục Ox. Khi vật qua vị trí cân bằng, tốc độ của nó là 20 cm/s. Khi vật ở biên, gia tốc của vật có độ lớn là 0,8 $m/s^2$. Khi vật cách vị trí cân bằng 4 cm thì nó có tốc độ
\choice
{\True 12 cm/s}
{20 cm/s}
{25 cm/s}
{25 cm/s}
\loigiai{
\begin{itemize}
\item Ta có: $v_{\max} = 20\ cm/s;\ a_{\max} = 0,8\ m/s^2 = 80\ cm/s^2 \Rightarrow \omega  = 4\ rad/s;\ A = 5\ cm$.
\item Khi $\left| x \right|=4\ cm$ thì $x^2+\dfrac{v^2}{\omega^2}=A^2\Rightarrow {{4}^2}+\dfrac{v^2}{{{4}^2}}={{5}^2}\Rightarrow \left| v \right|=12$ cm/s.
\end{itemize}
}

%\haidong{5}
\end{ex} \haidong{20}
%\loai{Dùng công thức độc lập x,v tìm $\omega$, T, f}
\begin{ex}%[3P1K2-2][Loai1]
Một vật dao động điều hoà với biên độ 4 cm. Khi nó có li độ là 2 cm thì vận tốc là 1 m/s. Tần số dao động là
\choice 
{1 Hz}
{3 Hz}
{1,2 Hz}
{\True 4,6 Hz}
\loigiai{
Áp dụng công thức độc lập với thời gian: $v^2 = \omega^2+(A^2-x^2)\Rightarrow \omega = 28,868\ rad/s \Rightarrow f = \dfrac{2\pi}{\omega} = 4,6\ Hz$.
}%<MyLT1>
%\haidong{4}
\end{ex} \haidong{20}
\begin{ex}%[3P1B2-2][Loai1]
Một vật dao động điều hòa với quỹ đạo dài 20 cm. Khi vật đi qua li độ 6 cm thì nó có tốc độ là $8\pi$  cm/s. Chu kì dao động của vật là
\choice
{4 s}
{0,5 s}
{\True 2 s}
{1 s}
\loigiai{
\begin{itemize}
\item Ta có: $A=\dfrac{L}{2}=10$ cm.
\item Luôn có: $x^2+\dfrac{v^2}{\omega^2}=A^2\Rightarrow 6^2+\dfrac{\left(8\pi \right)^2}{\omega^2}=10^2\Rightarrow \omega =\pi\ rad/s \Rightarrow T = 2\ s$.
\end{itemize}
}
%\haidong{4}
\end{ex} \haidong{20}
%\loai{Dùng công thức độc lập x,v tìm $v_{max}$}
\begin{ex}
Một chất điểm dao động điều hoà. Biết li độ và vận tốc của chất điểm tại thời điểm $t_1$ lần lượt là $x_1=3\ cm$ và $v_1=-60\sqrt 3\ cm/s$; tại thời điểm $t_2$ lần lượt là $x_2=3\sqrt 2\ cm$ và $v_2=60\sqrt 2\ cm/s$. Biên độ và tần số góc của dao động lần lượt bằng
\choice
{$6\ cm;\ 10\ rad/s$}
{\True $6\ cm;\ 20\ rad/s$}
{$12\ cm;\ 20\ rad/s$}
{$12\ cm;\ 10\ rad/s$}
\loigiai{
\begin{itemize}
\item Áp dụng công thức độc lập thời gian.

\end{itemize}
}
%\haidong{6}
\end{ex} \haidong{20}
\begin{ex}
Một vật dao động điều hoà với tần số góc $4\ rad/s$. Khi vật đi qua li độ $-8 \ cm$ thì nó có tốc độ là $32 \ cm/s$. Biên độ dao động của vật bằng
\choice
{$8,0\ cm$}
{$10,4\ cm$}
{\True $11,3\ cm$}
{$14,0\ cm$}
\loigiai{
$11,3\ cm$.
}
\end{ex}
\haidong{20}
\begin{ex}
Một vật dao động điều hòa với tần số góc là $\omega = 10 \ rad/s$. Khi vật cách vị trí cân bằng $9 \ cm$ thì tốc độ của vật là $120 \ cm/s$. Biên độ dao động của vật bằng
\choice
{$10\ cm$}
{$12,5\ cm$}
{\True $15\ cm$}
{$18\ cm$}
\loigiai{
}
\end{ex}
\begin{ex}
Một dao động điều hòa có vận tốc và li độ tại thời điểm $t_{1}$ và $t_{2}$ tương ứng là  $v_{1}=20 \ cm/s$; $x_{1}=8 \sqrt{3} \ cm$ và $v_{2}=20 \sqrt{2} \ cm/s$; $x_{2}=8 \sqrt{2} \ cm$. Vận tốc của vật có độ lớn cực đại bằng
\choice
{$25\ cm/s$}
{$32\ cm/s$}
{\True $40\ cm/s$}
{$50\ cm/s$}
\loigiai{
Ta có mối liên hệ giữa vận tốc, li độ và biên độ trong dao động điều hòa: $v^2 = \omega^2 (A^2 - x^2)$.\\
Áp dụng cho hai thời điểm $t_1$ và $t_2$:\\
$$v_1^2 = \omega^2 (A^2 - x_1^2) \implies (20)^2 = \omega^2 (A^2 - (8\sqrt{3})^2) \implies 400 = \omega^2 (A^2 - 192) \quad (1)$$
$$v_2^2 = \omega^2 (A^2 - x_2^2) \implies (20\sqrt{2})^2 = \omega^2 (A^2 - (8\sqrt{2})^2) \implies 800 = \omega^2 (A^2 - 128) \quad (2)$$
Chia (2) cho (1):\\
$$\dfrac{800}{400} = \dfrac{\omega^2 (A^2 - 128)}{\omega^2 (A^2 - 192)} \Rightarrow 2 = \dfrac{A^2 - 128}{A^2 - 192}$$
$$2(A^2 - 192) = A^2 - 128 \Rightarrow A^2 = 256 \Rightarrow A = 16\ cm$$
Thay vào (1):\\
$$400 = \omega^2 (256 - 192) = 64\omega^2 \Rightarrow \omega^2 = \dfrac{25}{4} \Rightarrow \omega = 2,5\ rad/s$$
Vận tốc cực đại: $v_{\max} = \omega A = 2,5 . 16 = 40\ cm/s$.
}
\end{ex}

\begin{ex}%[3P1B2-2][Loai1]
Vật dao động điều hòa. Khi vật có li độ 3 cm thì tốc độ của nó là $15\sqrt{3}$ cm/s, khi nó có li độ $3\sqrt{2}$ cm thì tốc độ của nó là $15\sqrt{2}$ cm/s. Tốc độ của vật khi đi qua vị trí cân bằng là
\choice
{50 cm/s}
{\True 30 cm/s}
{25 cm/s}
{20 cm/s}
\loigiai{
\begin{itemize}
\item $\begin{cases} 
x_1^2 + \dfrac{v_1^2}{\omega^2}=A^2\quad (1)\\
x_2^2 + \dfrac{v_2^2}{\omega^2} = A^2\quad (2)
\end{cases} 
\Rightarrow x_1^2 + \dfrac{v_1^2}{\omega^2}=x_2^2+\dfrac{v_2^2}{\omega^2} \Rightarrow \omega = \sqrt{\dfrac{v_1^2-v_2^2}{x_2^2-x_1^2}}=5\ rad/s$. 
\item $\left( 1 \right)\Rightarrow A^2=x_1^2+\dfrac{v_1^2}{\omega^2} \Rightarrow A=6\ cm\Rightarrow v_{\max} = A\omega  = 30$ cm/s.
\end{itemize}
}
%\haidong{6}
\end{ex} \haidong{20}

%\loai{Công thức độc lập giữa a và x}
\begin{ex} %[3P1B2-2][Loai1]  
Cho một chất điểm dao động điều hòa với tần số 2 Hz. Tốc độ cực đại trong quá trình dao động bằng 16$\pi$ cm/s. Khi ngang qua vị trí có li độ 2 cm, gia tốc chuyển động của chất điểm có độ lớn bằng
\choice
{$8\pi^2\ cm/s^2$}
{$64\pi^2\ cm/s^2$}
{$8\pi^2\sqrt{3}\ cm/s^2$}
{\True $32\pi^2\ cm/s^2$}
\loigiai{
\begin{itemize}
\item Ta có: $f = 2\ Hz \Rightarrow \omega = 4\pi$ rad/s.
\item Do x và a ngược pha nên $a = -\left(4\pi\right)^2.2 = -32\pi^2\ cm/s^2$.
\item Vậy gia tốc chuyển động của chất điểm có độ lớn bằng $32\pi^2\ cm/s^2.$
\end{itemize}
}

%\haidong{8}
\end{ex} \haidong{20}
%\loai{Cho hai gia tốc. Tìm gia tốc tại vị trí thứ 3}
\begin{ex} %[3P1-2.2K][Loai1]  
Một chất điểm đang dao động điều hòa trên một đoạn thẳng. Gia tốc của chất điểm khi ngang qua vị trí M và N lần lượt là $a_M = 30\ cm/s^2$ và $a_N = 70\ cm/s^2$. Khi đi ngang qua trung điểm I của đoạn MN, gia tốc chuyển động của chất điểm là
\choice
{$30\ cm/s^2$}
{$40\ cm/s^2$}
{\True $50\ cm/s^2$}
{$26\ cm/s^2$}
\loigiai{
\begin{itemize}
\item Ta có: $\begin{cases}
a_M = -\omega^2.x_M = 30\ cm/s^2.\\ 
a_N = -\omega^2.x_N = 70\ cm/s^2.\end{cases}$
\item I là trung điểm của đoạn MN thì $x_I = \dfrac{x_M + x_N}{2}$\\
$\Rightarrow a_I = -\omega^2.\dfrac{x_M + x_N}{2} = \dfrac{a_M + a_N}{2} = \dfrac{30 + 70}{2} = 50\ cm/s^2$.
\end{itemize}}
%\haidong{12}
\end{ex} \haidong{20}
\begin{ex}
Một chất điểm dao động điều hoà có hệ thức liên hệ giữa li độ $x$ và vận tốc $v$ là $\dfrac{x^2}{16}+\dfrac{v^2}{640}=1$ (trong đó $x$ tính bằng $cm$, $v$ tính bằng $cm/s)$. Lấy $\pi^2=10$. Chu kì dao động của chất điểm bằng
\choice
{$0,5\ s$}
{\True $1\ s$}
{$1,5\ s$}
{$2\ s$}
\loigiai{
$1 \ s$
}
%\haidong{5}
\end{ex} \haidong{20}
\begin{ex}
Một chất điểm dao động điều hoà với phương trình $x = 10 \cos 2 t\ cm$ ($t$ tính bằng $s$). Tại thời điểm $t$, chất điểm qua vị trí cách vị trí cân bằng một đoạn $5\ cm$ thì chất điểm có tốc độ bằng
\choice
{$8,7\ cm/s$}
{\True $17,3\ cm/s$}
{$20,0\ cm/s$}
{$25,9\ cm/s$}
\loigiai{
$17,3\ cm/s$.
}
\end{ex}
\begin{ex}
Một chất điểm dao động điều hoà trên $O x$ với quỹ đạo dài $16 \ cm$. Tại thời điểm $t_0$, vận tốc và gia tốc của chất điểm lần lượt là $40 \ cm/s$ và $4 \sqrt{3} \ m/s^2$. Tần số góc của chất điểm là
\choice
{$5\ rad/s$}
{\True $10\ rad/s$}
{$12,6\ rad/s$}
{$17,3\ rad/s$}
\loigiai{
$10\ rad/s$
}
\end{ex}
\haidong{20}
\begin{ex}
Một chất điểm dao động điều hoà với phương trình li độ $x=5 \cos (\pi t+\dfrac{\pi}{3})\ cm$ ($t$ tính bằng giây). Lấy $\pi^2=10$. Gia tốc của chất điểm khi chất điểm qua vị trí có có li độ $3 \ cm$ có độ lớn bằng bao nhiêu $cm/s^2$.\\
\choice
{$15\ cm/s^2$}
{$25\ cm/s^2$}
{\True $30\ cm/s^2$}
{$50\ cm/s^2$}
\loigiai{
$30\ cm/s^2$
}
%\haidong{5}
\end{ex}
\haidong{20}

\begin{ex}
Một vật dao động điều hòa trên quỹ đạo dài $40\ cm$. Khi li độ của vật là $10\ cm$, vật có tốc độ $20\pi\sqrt{3}\ cm/s$. Chu kì dao động của vật là
\choice
{$0,5\ s$}
{\True $1\ s$}
{$1,5\ s$}
{$2\ s$}
\loigiai{
Biên độ $A = \dfrac{40}{2}=20\ cm$.\\
$$\frac{v^{2}}{\omega^{2}}+x^{2}=A^{2}\Rightarrow 
\frac{\left(20\pi\sqrt{3}\right)^{2}}{\omega^{2}}+10^{2}=20^{2}
\Rightarrow \omega=2\pi\ rad/s$$
$$T=\dfrac{2\pi}{\omega}=1\ s.$$
}
\end{ex}
\setcounter{ex}{0}
\PhanII
\begin{ex}
Trong dao động điều hoà, xét các nhận định sau đây về mối liên hệ giữa li độ $x$, vận tốc $v$, vận tốc cực đại $v_{\max}$, tần số góc $\omega$ và chu kì $T$.
\choiceTF[t]
{\True Luôn tồn tại hệ thức $\displaystyle \left(\frac{x}{A}\right)^{2} + \left(\frac{v}{v_{\max}}\right)^{2} = 1$}
{$v_{\max}$ được tính bằng $v_{\max}=A\omega^{2}$}
{\True Chu kì dao động thoả mãn $T = \dfrac{2\pi}{\omega}$}
{Khi vật ở vị trí biên, gia tốc của vật bằng $0$}
\loigiai{
\begin{itemchoice}
\itemch
\itemch
\itemch
\itemch
\end{itemchoice}
}
\end{ex}

\begin{ex}
Một chất điểm dao động điều hoà với phương trình vận tốc $v=10 \pi \cos (\pi t+\dfrac{\pi}{3})\ cm/s$ ($t$ tính bằng giây). Lấy $\pi^2$ $=10$.
\choiceTF[t]
{\True Quỹ đạo dao động của vật là một đoạn thẳng dài $20 \ cm$}
{Tại thời điểm ban đầu vật đang đi theo chiều âm}
{\True Gia tốc cực đại của vật bằng $1 \ m/s^2$}
{\True Khi chất điểm có vận tốc $5 \pi\ cm/s$ thì gia tốc của chất điểm có độ lớn bằng $\dfrac{\sqrt{3}}{2} \ m/s^2$}
\loigiai{
\begin{itemchoice}
\itemch
\itemch
\itemch
\itemch
\end{itemchoice}
}
%\haidong{5}
\end{ex} \haidong{20}
\begin{ex}
Một vật dao động điều hòa có phương trình: $x = 5 \cos(4\pi t - \dfrac{\pi}{6})$ trong đó x tính bằng $cm$, $t$ tính bằng giây.
\choiceTF[t]
{\True Biên độ dao động của vật là $5\ cm$}
{Chu kì dao động của vật là $2\ s$}
{\True Gia tốc của vật tại vị trí có li độ $3\ cm$ bằng $-48\pi^{2}\ cm/s^2$}
{Pha dao động của vật tại thời điểm ban đầu bằng $0,6\pi\ rad$}
\loigiai{
\begin{itemchoice}
\itemch
\itemch
\itemch
\itemch
\end{itemchoice}
}
\end{ex}
\begin{ex}
Cho phương trình dao động điều hòa của một vật có dạng: $x=6 \cos \left(4 \pi t-\dfrac{\pi}{3}\right)\ cm$.
\choiceTF[t]
{Khi pha dao động bằng $\dfrac{2\pi}{3}\ rad$ thì gia tốc của vật bằng $80\pi^2\ cm/s^2$}
{Chiều dài quỹ đạo của vật là $10\ cm$}
{Tốc độ cực đại của vật bằng $40\pi\ cm/s$}
{\True Chu kì dao động của vật bằng $0,5\ s$}
\loigiai{ }
\end{ex}

\setcounter{ex}{0}
\PhanIII
\begin{ex}
Một chất điểm dao động điều hoà có chu kì $2 \ s$ và biên độ $10 \ cm$. Khi chất điểm qua trị trí cách vị trí biên $5 \ cm$ thì có tốc độ là bao nhiêu $cm/s$ (làm tròn kết quả đến chữ số hàng phần mười)?
\shortans[oly]{27,2}
\loigiai{
$27,2 \ cm/s$
}
%\haidong{4}
\end{ex} \haidong{20}
\begin{ex}
Một chất điểm dao động điều hoà với phương trình vận tốc $v=10 \pi \cos (\pi t+\dfrac{\pi}{3})\ cm/s$ ($t$ tính bằng giây). Lấy $\pi^2=10$. Khi chất điểm có li độ là $5 \sqrt{3} \ cm$ thì tốc độ của chất điểm bằng bao nhiêu $cm/s$ (làm tròn kết quả đến chữ số hàng phần mười)? 
\shortans[oly]{15,8}
\loigiai{
}
%\haidong{5}
\end{ex} \haidong{20}

\begin{ex}
Một chất điểm đang dao dộng điều hoà dọc theo trục $O x$ với phương trình li độ là $x=4 \cos (10 \pi t+\varphi)\ cm$. Khi góc pha dao động là $\dfrac{7 \pi}{6}$ rad thì chất điểm đó có vận tốc bằng bao nhiêu $cm/s$ (làm tròn kết quả đến chữ số hàng phần trăm)?
\shortans[oly]{62,8}
\loigiai{
$62,8 \ cm/s$
}
%\haidong{3}
\end{ex} \haidong{20}


\begin{ex}
Một chất điểm dao động điều hoà dọc theo trục $O x$. Tại thời điểm $t_1$, li độ của chất điểm bằng $3 \ cm$ và vận tốc bằng $6 \pi \sqrt{3} \ cm/s$. Tại thời điểm $t_2$, li độ bằng $3 \sqrt{2} \ cm$ và vận tốc bằng $6 \pi \sqrt{2} \ cm/s$. Tốc độ cực đại của chất điểm bằng bao nhiêu $cm/s$ (làm tròn kết quả đến chữ số hàng phần mười)?
\shortans[oly]{37,7}
\loigiai{
$37,7 \ cm/s$
}
%\haidong{7}
\end{ex} \haidong{20}

\begin{ex}
Một chất điểm dao động điều hoà dọc theo trục $O x$ với biên độ $5 \ cm$ và chu kì $2 \ s$. Lấy $\pi^2=10$. Tại thời điểm $t_0$, chất điểm có tốc độ $2,5 \pi\ cm/s$ thì độ lớn gia tốc của chất điểm là bao nhiêu $\ cm/s^2$ (làm tròn kết quả đến chữ số hàng phần mười)?
\shortans[oly]{43,3} \loigiai{
$43,3\ cm/s^2$
}
%\haidong{5}
\end{ex} \haidong{20}

\begin{ex}
Một chất điểm dao động điều hoà có hệ thức liên hệ giữa vận tốc $v$ và gia tốc $a$ là $\dfrac{v^2}{640}+\dfrac{a^2}{2,56}=1$ (trong đó, $v$ tính bằng $cm/s$ và $a$ tính bằng $m/s^2$). Biên độ dao động của chất điểm bằng bao nhiêu $cm$?
\shortans[oly]{1}
\loigiai{
$4 \ cm$}
%\haidong{4}
\end{ex} \haidong{20}